\section{A minimal counterexample and neighborhood expansion}
\label{sec:reduction}

All graphs in the proof are finite and simple. For a graph $H$, write
$\alpha(H)$ for its independence number, $\omega(H)$ for its clique
number, and $\delta(H)$ for its minimum degree. If $X\subseteq V(H)$,
then $H[X]$ is the induced subgraph on $X$ and $H-X$ is the induced
subgraph on $V(H)\setminus X$. We sometimes write $\alpha(X)$ for
$\alpha(H[X])$ when the ambient graph is clear. The open neighborhood
of a vertex $x$ is denoted by $N_H(x)$. For a set $X$, put
$N_H(X)=\bigcup_{x\in X}N_H(x)$, and
\[
  N_H[X]=X\cup\bigcup_{x\in X}N_H(x)
\]
is the closed neighborhood of a set $X$. Thus $N_H[x]=N_H[\{x\}]$.
The order of a path or cycle is its number of vertices, and its length
is its number of edges. In particular, a copy of $C_m$ always has
exactly $m$ vertices.

The upper bound is equivalent to the assertion that every graph on
$(m-1)(n-1)+1$ vertices contains either a cycle of order $m$ or an
independent set of order $n$: apply this assertion to the red graph of
a two-coloring. Set $k=m-1$ and $a=n-1$. In the range $m\ge n\ge3$
with $(m,n)=(3,3)$ excluded, these parameters satisfy
\[
  k\ge3,\qquad 2\le a\le k.
\]
Conversely, every such pair $(k,a)$ corresponds to a pair in the
required range. We will rule out counterexamples for each fixed $k$.

Fix $k\ge3$ and suppose that the upper bound fails for at least one
$a\in\{2,\ldots,k\}$. Choose the least such $a$, and fix a graph $G$
with
\begin{equation}
  |V(G)|=ka+1,\qquad \alpha(G)\le a,
  \qquad G\text{ contains no }C_{k+1}.
  \label{red:counterexample}
\end{equation}
The choice of the least independence parameter will turn deletion of
an independent set and all its neighbors into a useful expansion
bound. The graph $G$ and the parameter $k$ retain this meaning throughout
the upper-bound proof; the parameter $a$ is used through
Section~\ref{sec:small}.

\begin{lemma}\label{red:smaller}
For every integer $b$ with $0\le b<a$, a graph $H$ containing no
$C_{k+1}$ and satisfying $\alpha(H)\le b$ has at most $kb$ vertices.
\end{lemma}

\begin{proof}
If $b=0$, the graph $H$ is empty, since any vertex would be an
independent set of order one. If $b=1$, the graph is complete. A
complete graph with at least $k+1$ vertices contains a cycle of exactly
$k+1$ vertices, so in this case $|V(H)|\le k$.

Now let $2\le b<a$. If $|V(H)|\ge kb+1$, take an induced subgraph on
$kb+1$ vertices. It still contains no $C_{k+1}$ and has independence
number at most $b$. It would therefore be a counterexample for the
smaller parameter $b$, contrary to the choice of $a$.
\end{proof}

Independent-set deletion and neighbourhood expansion already enter the
cycle--clique argument of Erd\H{o}s, Faudree, Rousseau and Schelp
\cite[Section~3]{ErdosFaudreeRousseauSchelp1978}. The following closed
neighbourhood bound comes directly from our fixed-parameter minimality.

\begin{lemma}[Independent-set expansion]\label{red:expansion}
For every nonempty independent set $I$ in the graph $G$ fixed in
\eqref{red:counterexample},
\begin{equation}
  |N_G[I]|\ge k|I|+1.
  \label{red:expansion-bound}
\end{equation}
In particular, $\delta(G)\ge k$ and $\alpha(G)\le k$.
\end{lemma}

\begin{proof}
Put $r=|I|$, so that $1\le r\le a$. Every independent set in
$G-N_G[I]$ can be adjoined to $I$, because it has no edges to $I$.
Consequently,
\[
  \alpha\bigl(G-N_G[I]\bigr)\le a-r.
\]
The integer $a-r$ lies between zero and $a-1$. The graph
$G-N_G[I]$ also contains no $C_{k+1}$, so Lemma~\ref{red:smaller}
applies, including when $r=a$, and gives
\[
  |V(G)\setminus N_G[I]|\le k(a-r).
\]
Subtracting from $|V(G)|=ka+1$ proves
\eqref{red:expansion-bound}. Applying the bound to $I=\{x\}$ yields
$d_G(x)+1\ge k+1$, and hence $\delta(G)\ge k$. Finally,
$\alpha(G)\le a\le k$ by the choice of the parameters.
\end{proof}

The remainder of the upper-bound proof uses
\eqref{red:counterexample} and Lemma~\ref{red:expansion} to derive a
contradiction. The first structural step, in
Section~\ref{sec:clique}, finds a large clique. Section~\ref{sec:small}
then disposes of $k=3,4$; all later arguments may assume $k\ge5$.
